Fix \(\mathcal M\in\mathfrak F(\mathcal D)\) with \(b\in\operatorname{col}(B_x)\). Theorem~\ref{thm:target-span-iff} and Definition~\ref{def:unconditional-weight-map} provide a fixed vector \(\kappa\in\mathbb R^{\lvert\mathcal E\rvert}\) such that Lemma~\ref{lem:observable-factorization} gives the two population restrictions
\[
H_x\kappa=b,
\qquad
z_{x,y}^{\top}\kappa=\theta_{x,y}.
\tag{1}
\]
No norm bound on \(\kappa\) is imposed.

We first derive the concentration inequality used below. If \(Z\) is Bernoulli with mean \(p_Z\), define
\[
\psi(s)=\log\mathbb E\!\left[\exp\{s(Z-p_Z)\}\right].
\]
Under exponential tilting, \(\psi''(s)=p_s(1-p_s)\), where \(p_s\in[0,1]\) is the tilted success probability, and therefore \(\psi''(s)\le1/4\). Since \(\psi(0)=\psi'(0)=0\), integrating this second-derivative bound twice gives \(\psi(s)\le s^2/8\) for every \(s\in\mathbb R\). Thus, for independent Bernoulli variables \(Z_1,\ldots,Z_N\) with common mean \(p_Z\), put \(T=\sum_{i=1}^N(Z_i-p_Z)\). On \(\{T\ge Na\}\), one has \(e^{sT}\ge e^{sNa}\), and therefore
\[
\mathbb E[e^{sT}]\ge e^{sNa}\Pr\{T\ge Na\}.
\]
Rearranging this inequality and then using independence gives, for \(a>0\) and \(s>0\),
\[
\begin{aligned}
\Pr\!\left\{N^{-1}\sum_{i=1}^N Z_i-p_Z\ge a\right\}
&\le \exp(-sNa)\prod_{i=1}^N
 \mathbb E\!\left[\exp\{s(Z_i-p_Z)\}\right]\\
&\le\exp\!\left(-sNa+\frac{Ns^2}{8}\right).
\end{aligned}
\tag{2}
\]
Taking \(s=4a\), and applying the same argument to \(-Z_i\), yields the self-contained two-sided bound
\[
\Pr\!\left\{\left|N^{-1}\sum_{i=1}^N Z_i-p_Z\right|>a\right\}
\le2e^{-2Na^2}.
\tag{3}
\]

Every coordinate of \(\widehat H_x\) and \(\widehat z_{x,y}\) is an average of source-sample Bernoulli indicators, with respective means \(H_x(w,e)\) and \(z_{x,y}(e)\), by Assumption~\ref{ass:source-iid-sampling} and Definition~\ref{def:unconditional-balancing-moments}. Every coordinate of \(\widehat b\) is an average of target-sample Bernoulli indicators with mean \(b_w\), by Assumption~\ref{ass:target-iid-sampling}. Assumption~\ref{ass:sample-source-independence} makes their joint law the stated product law, although the union bound below does not require independence across coordinates. Define the simultaneous event
\[
\begin{aligned}
\mathcal A=\{&\max_{w,e}|\widehat H_x(w,e)-H_x(w,e)|\le r_s,\\
&\max_e|\widehat z_{x,y}(e)-z_{x,y}(e)|\le r_s,\\
&\max_w|\widehat b_w-b_w|\le r_t\}.
\end{aligned}
\tag{4}
\]
There are exactly \(L=m\lvert\mathcal E\rvert+\lvert\mathcal E\rvert+m\) coordinates in (4). From the definitions of \(r_s\) and \(r_t\), equation (3) bounds the failure probability of each coordinate by
\[
2\exp(-2n_s r_s^2)=\frac{\alpha}{L}
\quad\text{or}\quad
2\exp(-2n_t r_t^2)=\frac{\alpha}{L},
\tag{5}
\]
as appropriate. The union bound therefore gives
\[
\mathbb P_{\mathcal M}^{\,n_s,n_t}(\mathcal A)\ge1-\alpha.
\tag{6}
\]

On \(\mathcal A\), use the same population balancing vector \(\kappa\) as in (1). For every \(w\in\mathcal W\),
\[
\begin{aligned}
| (\widehat H_x\kappa-\widehat b)_w |
&\le |((\widehat H_x-H_x)\kappa)_w|+|b_w-\widehat b_w|\\
&\le r_s\sum_{e\in\mathcal E}|\kappa(e)|+r_t
=r_s\lVert\kappa\rVert_1+r_t,
\end{aligned}
\tag{7}
\]
where (1) was used in the first line. Similarly,
\[
|\widehat z_{x,y}^{\top}\kappa-\theta_{x,y}|
=| (\widehat z_{x,y}-z_{x,y})^{\top}\kappa |
\le r_s\lVert\kappa\rVert_1.
\tag{8}
\]
Equations (7)--(8) are precisely the two feasibility inequalities in Definition~\ref{def:concentration-projection-set} with candidate \(\vartheta=\theta_{x,y}\). Since \(\theta_{x,y}\in[0,1]\), they imply
\[
\mathcal A\subseteq
\{\theta_{x,y}\in\widehat C_{x,y,1-\alpha}\}.
\tag{9}
\]
Combining (6) and (9) proves the advertised coverage.

All constants in (3)--(6) are numerical and the feasibility bounds retain the actual factor \(\lVert\kappa\rVert_1\). No lower bound on a singular value, no rank estimate, no covariance inverse, and no upper bound on the balancing norm entered the argument. Therefore the same inequality holds uniformly over every rank, every nonzero-singular-value configuration, every positive observed-cell law, and every magnitude of a balancing vector. The definition itself permits the realized projection set to be empty on \(\mathcal A^c\), or to equal \([0,1]\); neither possibility changes (9).